\documentclass[12pt, a4paper]{article}
% --- 1. Cấu hình Tiếng Việt và Font chữ chuẩn ---
\usepackage[utf8]{inputenc}
\usepackage[T5]{fontenc}
\usepackage[vietnamese]{babel}
\usepackage{mathptmx} % Font Times New Roman đồng bộ cả chữ và công thức toán, không gây xung đột gói

\makeatletter
\renewcommand{\normalsize}{\@setfontsize\normalsize{13pt}{16pt}}
\makeatother
\normalsize

% --- 2. Gói Toán học & Ký hiệu bổ trợ ---
\usepackage{amsmath, amssymb, amsfonts, mathrsfs, amsthm}
\usepackage{graphicx}
\usepackage{fontawesome}
\usepackage{booktabs, tabularx, array, multirow}
\usepackage{enumitem}

% --- 3. Đồ họa TikZ, Bảng biến thiên & Hình không gian ---
\usepackage{tikz}
\usetikzlibrary{calc, intersections, angles, quotes, patterns, arrows.meta, 3d}
\usepackage{pgfplots}
\pgfplotsset{compat=1.18}
\usepackage{tkz-euclide}
\usepackage{tkz-tab}

\renewcommand{\vec}[1]{\overrightarrow{#1}}

% --- 4. Hộp sư phạm (Tcolorbox) ---
\usepackage[many]{tcolorbox}
\tcbuselibrary{skins, breakable}

% --- 5. Căn lề chuẩn văn bản giáo án ---
\usepackage{geometry}
\geometry{
    a4paper,
    left=25mm,
    right=15mm,
    top=20mm,
    bottom=20mm
}

% --- 6. Header / Footer chuẩn hóa ---
\usepackage{fancyhdr}
\usepackage{lastpage}
\pagestyle{fancy}
\fancyhf{}

\fancyhead[L]{\small \textit{Kế hoạch bài dạy Toán 11}}
\fancyhead[R]{\textbf{Trang \thepage/\pageref{LastPage}}}
\fancyfoot[L]{\small \textit{Tổ Toán - Tin}}
\fancyfoot[R]{\small \textit{Trường THCS\&THPT Hồng Vân}}
\renewcommand{\headrulewidth}{0.4pt}
\renewcommand{\footrulewidth}{0.4pt}

% --- 7. Giãn dòng & Giãn đoạn theo chuẩn Nghị định 30/2020/NĐ-CP ---
\usepackage{setspace}
\setstretch{1.15}
\setlength{\parindent}{1.0cm}
\setlength{\parskip}{5pt}

% Định nghĩa hộp sư phạm chuyên dụng
\newtcolorbox{pedabox}[2][]{
    enhanced,
    breakable,
    colback=blue!3!white,
    colframe=blue!65!black,
    fonttitle=\bfseries\small,
    title={#2},
    arc=2mm,
    boxrule=0.6pt,
    left=3mm, right=3mm, top=2mm, bottom=2mm,
    #1
}

\newtcolorbox{aibox}[2][]{
    enhanced,
    breakable,
    colback=teal!4!white,
    colframe=teal!70!black,
    fonttitle=\bfseries\small,
    title={#2},
    arc=2mm,
    boxrule=0.6pt,
    left=3mm, right=3mm, top=2mm, bottom=2mm,
    #1
}

\newtcolorbox{scaffoldbox}[2][]{
    enhanced,
    breakable,
    colback=orange!4!white,
    colframe=orange!80!black,
    fonttitle=\bfseries\small,
    title={#2},
    arc=2mm,
    boxrule=0.6pt,
    left=3mm, right=3mm, top=2mm, bottom=2mm,
    #1
}

\newtcolorbox{stembox}[2][]{
    enhanced,
    breakable,
    colback=purple!4!white,
    colframe=purple!75!black,
    fonttitle=\bfseries\small,
    title={#2},
    arc=2mm,
    boxrule=0.6pt,
    left=3mm, right=3mm, top=2mm, bottom=2mm,
    #1
}

\begin{document}

\begin{center}
    {\fontsize{14pt}{17pt}\selectfont \textbf{KẾ HOẠCH BÀI DẠY MÔN TOÁN LỚP 11}}\\[6pt]
    {\fontsize{16pt}{19pt}\selectfont \textbf{BÀI 28. BIẾN CỐ HỢP, BIẾN CỐ GIAO, BIẾN CỐ ĐỘC LẬP}}\\[4pt]
    {\fontsize{12pt}{15pt}\selectfont \textbf{Môn học: Toán 11} \ \textbar \ \textbf{Thời lượng: 3 tiết (Tiết 83, 84, 85 theo PPCT | Tuần 30)}}
\end{center}

\vspace{5pt}

\section*{I. MỤC TIÊU}

\subsection*{1. Kiến thức, Năng lực}

\begin{itemize}[leftmargin=1.2cm]
    \item \textbf{Yêu cầu cần đạt (Nguyên văn CT GDPT 2018 Toán 11):}
    \begin{itemize}[leftmargin=0.6cm]
        \item Nhận biết được một số khái niệm về xác suất cổ điển: hợp và giao của các biến cố; biến cố độc lập.
        \item Mô tả được các biến cố hợp, biến cố giao bằng ngôn ngữ tập hợp và sơ đồ Ven.
        \item Phân biệt được biến cố hợp ("$A$ hoặc $B$"), biến cố giao ("cả $A$ và $B$"), biến cố xung khắc và biến cố độc lập trong các tình huống thực tiễn.
    \end{itemize}

    \item \textbf{Năng lực Toán học đặc thù:}
    \begin{itemize}[leftmargin=0.6cm]
        \item \textit{Năng lực tư duy và lập luận toán học:} So sánh, phân tích sự tương đồng giữa phép toán tập hợp (hợp, giao, phần bù) với các phép toán biến cố trong xác suất; phân biệt ranh giới bản chất giữa hai khái niệm dễ nhầm lẫn: \textit{biến cố xung khắc} (quan hệ không thể cùng xảy ra: $A \cap B = \emptyset$) và \textit{biến cố độc lập} (quan hệ về mặt thông tin, sự xảy ra của biến cố này không làm thay đổi xác suất của biến cố kia).
        \item \textit{Năng lực mô hình hóa toán học:} Mô tả các hiện tượng ngẫu nhiên trong đời sống và kỹ thuật (gieo xúc xắc, rút bài tây, bắn súng, kiểm tra chất lượng linh kiện điện tử, sàng lọc bệnh y tế) bằng các biến cố ngẫu nhiên; thiết lập không gian mẫu $\Omega$ và xác định chính xác tập hợp các kết quả thuận lợi cho biến cố hợp $A \cup B$ và biến cố giao $AB$.
        \item \textit{Năng lực giải quyết vấn đề toán học:} Nhận diện chính xác các từ khóa trong đề bài: "hoặc", "ít nhất một trong hai" $\to$ Biến cố hợp; "và", "đồng thời", "cả hai" $\to$ Biến cố giao; xác định tính độc lập của hai biến cố trong các phép thử lặp hoặc độc lập nhau.
        \item \textit{Năng lực giao tiếp toán học:} Trình bày lập luận xác suất rõ ràng, khoa học; sử dụng đúng chuẩn các ký hiệu toán học: $A \cup B, A \cap B, AB, \emptyset, P(A), \Omega$; diễn đạt chính xác mệnh đề bằng lời sang ngôn ngữ biến cố và ngược lại.
        \item \textit{Năng lực sử dụng công cụ, phương tiện học toán:} Sử dụng trực quan bộ dụng cụ thực hành xác suất (con xúc xắc, bộ bài 52 lá, túi bi màu); ứng dụng phần mềm vẽ sơ đồ số (Canva, GeoGebra, PowerPoint) để vẽ sơ đồ Ven thể hiện mối quan hệ giữa các biến cố.
    \end{itemize}

    \item \textbf{Năng lực số (Theo Thông tư 02/2025/TT-BGDĐT):}
    \begin{itemize}[leftmargin=0.6cm]
        \item \textit{Mã 3.1NC1a (Tạo và biểu diễn nội dung số trực quan):} Học sinh sử dụng phần mềm đồ họa hoặc sơ đồ tư duy số (như Canva, GeoGebra hoặc Google Slides) để tự tay thiết kế sơ đồ Ven minh họa không gian mẫu $\Omega$, tô màu trực quan phân biệt vùng biến cố hợp $A \cup B$, vùng biến cố giao $A \cap B$, vùng $A \setminus B$ và vùng biến cố xung khắc $A \cap B = \emptyset$.
    \end{itemize}

    \item \textbf{Năng lực AI (Theo Quyết định 2422/QĐ-BGDĐT):}
    \begin{itemize}[leftmargin=0.6cm]
        \item \textit{Mã 11.C5.2 (Khám phá mô hình xác suất trong hệ thống gợi ý AI):} Học sinh tìm hiểu cách các thuật toán Trí tuệ nhân tạo (Recommender Systems trên YouTube, TikTok, Google Search) mô hình hóa các sự kiện tìm kiếm và xem video của người dùng: phân tích các biến cố xem video độc lập hay phụ thuộc lẫn nhau để dự đoán xác suất người dùng nhấp chuột vào video tiếp theo, từ đó tối ưu hóa nội dung đề xuất.
    \end{itemize}

    \item \textbf{Năng lực chung:}
    \begin{itemize}[leftmargin=0.6cm]
        \item \textit{Tự chủ và tự học:} Chủ động chuẩn bị và thao tác với các đồ dùng ngẫu nhiên; tự giác suy nghĩ và giải quyết bài tập trên Phiếu học tập.
        \item \textit{Giao tiếp và hợp tác:} Hoạt động nhóm tích cực, phối hợp tung xúc xắc, rút thẻ bài và ghi chép thống kê kết quả thực nghiệm một cách trung thực.
    \end{itemize}
\end{itemize}

\subsection*{2. Phẩm chất}

\begin{itemize}[leftmargin=1.2cm]
    \item \textbf{Trung thực:} Tôn trọng tính khách quan của phép thử ngẫu nhiên; ghi chép chính xác kết quả gieo xúc xắc/rút bài thực tế, không tự ý bịa đặt số liệu.
    \item \textbf{Chăm chỉ:} Kiên nhẫn liệt kê đầy đủ các phần tử của không gian mẫu và tập hợp kết quả thuận lợi; cẩn thận kiểm tra từng trường hợp.
    \item \textbf{Trách nhiệm:} Hợp tác tích cực với bạn bè trong nhóm; giữ gìn vệ sinh và bảo quản tốt đồ dùng thực hành xác suất của nhà trường.
\end{itemize}

\section*{II. THIẾT BỊ DẠY HỌC VÀ HỌC LIỆU}

\begin{itemize}[leftmargin=1.2cm]
    \item \textbf{Giáo viên:}
    \begin{itemize}[leftmargin=0.6cm]
        \item Kế hoạch bài dạy chi tiết 3 tiết theo chuẩn Công văn 5512/BGDĐT-GDTrH.
        \item Bộ học cụ xác suất của trường: 8 hộp xúc xắc (mỗi hộp 2 con xúc xắc xanh - đỏ), 8 bộ bài tây 52 lá, túi đựng các quả cầu đánh số từ 1 đến 10.
        \item Slide bài giảng điện tử minh họa sinh động sơ đồ Ven, các hình ảnh trực quan về bộ bài, xúc xắc và thuật toán gợi ý video của AI.
        \item Hệ thống 3 Phiếu học tập phân tầng bậc thang (PHT số 1: Biến cố hợp; PHT số 2: Biến cố giao & Biến cố xung khắc; PHT số 3: Biến cố độc lập & AI).
        \item Bảng Rubric đánh giá năng lực học sinh.
    \end{itemize}

    \item \textbf{Học sinh:}
    \begin{itemize}[leftmargin=0.6cm]
        \item Sách giáo khoa Toán 11, vở ghi bài, bút màu (đỏ, xanh, vàng) để vẽ sơ đồ Ven.
        \item Mỗi bàn chuẩn bị 1 con xúc xắc hoặc 1 bộ bài tây (theo phân công từ tiết trước).
    \end{itemize}
\end{itemize}

\section*{III. TIẾN TRÌNH DẠY HỌC}

\begin{center}
    \textbf{\large TIẾT 1 (TIẾT 83 THEO PPCT | TUẦN 30)}\\[4pt]
    \textbf{\large BIẾN CỐ HỢP VÀ BIỂU DIỄN BẰNG SƠ ĐỒ VEN}
\end{center}

\subsection*{1. Hoạt động 1: Mở đầu (Khởi động - 6 phút)}

\begin{itemize}[leftmargin=1.2cm]
    \item[a)] \textbf{Mục tiêu:}
    \begin{itemize}[leftmargin=0.6cm]
        \item Tạo hứng thú học tập thông qua trò chơi rút bài tây quen thuộc; dẫn dắt học sinh đến khái niệm biến cố xuất hiện khi có từ "hoặc".
    \end{itemize}

    \item[b)] \textbf{Nội dung:}
    \begin{itemize}[leftmargin=0.6cm]
        \item Giáo viên cho một bộ bài tây 52 lá đã xáo đều. Rút ngẫu nhiên 1 lá bài.
        \item Xét hai biến cố:
        \begin{itemize}
            \item $A$: "Lá bài rút ra là lá Át (A)".
            \item $B$: "Lá bài rút ra mang chất Cơ (trái tim màu đỏ)".
        \end{itemize}
        \item Đặt câu hỏi: \textit{"Biến cố $C$: 'Lá bài rút ra là lá Át hoặc mang chất Cơ' sẽ xảy ra khi nào? Có bao nhiêu lá bài trong bộ bài làm cho biến cố $C$ xảy ra?"}
    \end{itemize}

    \item[c)] \textbf{Sản phẩm:}
    \begin{itemize}[leftmargin=0.6cm]
        \item Học sinh trả lời: Biến cố $C$ xảy ra khi rút được bất kỳ lá Át nào (Át Bích, Át Chuồn, Át Rô, Át Cơ) HOẶC bất kỳ lá chất Cơ nào (từ lá 2 Cơ đến K Cơ).
        \item Số lá bài thuận lợi: Có 4 lá Át, có 13 lá Cơ. Nhưng lá "Át Cơ" vừa là lá Át vừa mang chất Cơ (được tính trùng 1 lần). Do đó số lá bài thỏa mãn là $4 + 13 - 1 = 16$ lá bài.
    \end{itemize}

    \item[d)] \textbf{Tổ chức thực hiện:} GV nhận xét câu trả lời, giới thiệu biến cố $C$ chính là "Biến cố hợp" của hai biến cố $A$ và $B$.
\end{itemize}

\subsection*{2. Hoạt động 2: Hình thành kiến thức mới (22 phút)}

\begin{itemize}[leftmargin=1.2cm]
    \item[a)] \textbf{Mục tiêu:}
    \begin{itemize}[leftmargin=0.6cm]
        \item Nắm vững định nghĩa biến cố hợp của hai biến cố.
        \item Biết cách mô tả biến cố hợp dưới dạng tập hợp và biểu diễn trực quan bằng sơ đồ Ven.
        \item Rèn luyện năng lực số thông qua vẽ sơ đồ Ven trên thiết bị số (Mã 3.1NC1a).
    \end{itemize}

    \item[b)] \textbf{Nội dung:}
    \begin{itemize}[leftmargin=0.6cm]
        \item \textbf{1. Định nghĩa biến cố hợp:}
        \begin{itemize}
            \item Cho hai biến cố $A$ và $B$ cùng liên quan đến một phép thử.
            \item Biến cố \textit{"$A$ hoặc $B$ xảy ra"}, kí hiệu là $A \cup B$, được gọi là \textbf{biến cố hợp} của $A$ và $B$.
            \item \textit{Quy tắc nhận diện ngôn ngữ:} Biến cố hợp $A \cup B$ xảy ra khi \textbf{có ít nhất một trong hai biến cố} $A$ hoặc $B$ xảy ra.
        \end{itemize}

        \item \textbf{2. Mô tả bằng tập hợp và Sơ đồ Ven:}
        \begin{itemize}
            \item Giả sử biến cố $A$ được mô tả bởi tập hợp $A \subset \Omega$, biến cố $B$ được mô tả bởi tập hợp $B \subset \Omega$.
            \item Khi đó, tập hợp các kết quả thuận lợi cho biến cố hợp chính là hợp của hai tập hợp $A$ và $B$:
            $$C = A \cup B = \{\omega \in \Omega \mid \omega \in A \text{ hoặc } \omega \in B\}$$
            \item \textit{Biểu diễn sơ đồ Ven:} Không gian mẫu $\Omega$ là hình chữ nhật bao quanh. Vùng biểu diễn biến cố hợp $A \cup B$ là toàn bộ diện tích thuộc hình tròn $A$ hoặc hình tròn $B$ (kể cả phần chung).
        \end{itemize}
    \end{itemize}

    \item[c)] \textbf{Sản phẩm:} Học sinh ghi chép định nghĩa vào vở, vẽ sơ đồ Ven minh họa biến cố hợp và tô màu phân biệt rõ ràng.
    \item[d)] \textbf{Tổ chức thực hiện:}
\end{itemize}

\begin{pedabox}{Tiến trình tổ chức Hoạt động 2 (Kiến thức mới - Tiết 1)}
    \textbf{Bước 1: Chuyển giao nhiệm vụ}
    \begin{itemize}[leftmargin=0.6cm]
        \item GV yêu cầu học sinh làm việc cá nhân: Đọc định nghĩa trong SGK và phát biểu lại bằng lời của mình về ý nghĩa của từ "hoặc".
        \item Yêu cầu vẽ sơ đồ Ven của $A \cup B$ vào vở bằng 2 màu bút khác nhau.
    \end{itemize}

    \textbf{Bước 2: Thực hiện nhiệm vụ có giàn giáo sư phạm}
    \begin{itemize}[leftmargin=0.6cm]
        \item Học sinh quan sát, vẽ 2 hình tròn giao nhau bên trong hình chữ nhật $\Omega$.
        \item GV hướng dẫn học sinh trung bình - yếu: Tô màu toàn bộ hình tròn $A$, sau đó tô tiếp hình tròn $B$; phần nào đã tô màu thì chính là thuộc biến cố hợp $A \cup B$.
    \end{itemize}

    \textbf{Bước 3: Báo cáo, thảo luận}
    \begin{itemize}[leftmargin=0.6cm]
        \item GV mời 1 học sinh lên bảng vẽ sơ đồ Ven. Cả lớp quan sát, nhận xét.
    \end{itemize}

    \textbf{Bước 4: Kết luận, chuẩn hóa}
    \begin{itemize}[leftmargin=0.6cm]
        \item GV chốt định nghĩa và nhấn mạnh từ khóa nhận diện: "hoặc", "ít nhất một trong hai".
    \end{itemize}
\end{pedabox}

\begin{aibox}{Tích hợp Năng lực số (Theo Thông tư 02/2025/TT-BGDĐT - Mã 3.1NC1a)}
    \textbf{Thực hành thiết kế sơ đồ Ven số trên phần mềm đồ họa trực quan:}
    \begin{itemize}[leftmargin=0.6cm]
        \item Học sinh mở ứng dụng trình chiếu (Google Slides / PowerPoint / Canva) trên điện thoại hoặc máy tính phòng thực hành.
        \item Sử dụng công cụ vẽ hình khối (\texttt{Shapes}): Chèn hình chữ nhật đại diện không gian mẫu $\Omega$; chèn 2 hình tròn đan xen nhau đại diện cho biến cố $A$ và $B$.
        \item Sử dụng tính năng hợp hình khối (\texttt{Shape Union}) hoặc chỉnh độ trong suốt (\texttt{Transparency}) để làm nổi bật miền biến cố hợp $A \cup B$.
        \item Xuất ảnh sơ đồ Ven và chia sẻ lên nhóm Zalo học tập của lớp.
    \end{itemize}
\end{aibox}

\subsection*{3. Hoạt động 3: Luyện tập (12 phút)}

\begin{itemize}[leftmargin=1.2cm]
    \item[a)] \textbf{Mục tiêu:} Áp dụng thành thạo định nghĩa biến cố hợp để liệt kê các kết quả thuận lợi trong bài toán gieo xúc xắc.
    \item[b)] \textbf{Nội dung:} Thực hiện bài toán trên Phiếu học tập số 1:
    \begin{quote}
        \textbf{Ví dụ 1:} Gieo một con xúc xắc cân đối và đồng chất một lần. Gọi:
        \begin{itemize}
            \item $A$ là biến cố: "Số chấm xuất hiện là số chẵn".
            \item $B$ là biến cố: "Số chấm xuất hiện là số nguyên tố".
        \end{itemize}
        \begin{enumerate}
            \item Hãy liệt kê các phần tử của không gian mẫu $\Omega$ và các tập hợp thuận lợi cho biến cố $A, B$.
            \item Hãy phát biểu biến cố hợp $A \cup B$ bằng lời và liệt kê các kết quả thuận lợi cho $A \cup B$.
        \end{enumerate}
    \end{quote}

    \item[c)] \textbf{Sản phẩm (Lời giải chi tiết):}
    \begin{itemize}[leftmargin=0.6cm]
        \item Không gian mẫu: $\Omega = \{1; 2; 3; 4; 5; 6\}$.
        \item Các kết quả thuận lợi cho biến cố $A$: $A = \{2; 4; 6\}$.
        \item Các kết quả thuận lợi cho biến cố $B$: $B = \{2; 3; 5\}$ (các số nguyên tố từ 1 đến 6).
        \item Phát biểu biến cố $A \cup B$ bằng lời: "Số chấm xuất hiện là số chẵn hoặc số nguyên tố" (hoặc "Số chấm xuất hiện có ít nhất một tính chất: là số chẵn hoặc là số nguyên tố").
        \item Tập hợp các kết quả thuận lợi cho $A \cup B$:
        $$A \cup B = \{2; 3; 4; 5; 6\}$$
        (Lưu ý: Phần tử $2$ xuất hiện ở cả hai tập nhưng trong tập hợp hợp chỉ viết một lần).
    \end{itemize}

    \item[d)] \textbf{Tổ chức thực hiện:}
    \begin{itemize}[leftmargin=0.6cm]
        \item Học sinh làm việc cá nhân trong 5 phút. GV hỗ trợ học sinh trung bình nhớ lại số 2 là số nguyên tố chẵn duy nhất.
        \item Gọi 1 học sinh lên bảng viết kết quả; cả lớp nhận xét, đánh giá.
    \end{itemize}
\end{itemize}

\subsection*{4. Hoạt động 4: Vận dụng (5 phút)}

\begin{itemize}[leftmargin=1.2cm]
    \item[a)] \textbf{Mục tiêu:} Mở rộng bài toán biến cố hợp cho trường hợp 3 biến cố.
    \item[b)] \textbf{Nội dung & Hướng dẫn về nhà:}
    \begin{itemize}[leftmargin=0.6cm]
        \item Cho ba biến cố $A, B, C$. Biến cố "Có ít nhất một trong ba biến cố xảy ra" được viết dưới dạng ký hiệu là $A \cup B \cup C$.
        \item Chuẩn bị nội dung Tiết 2: Thế nào là biến cố giao ("đồng thời") và thế nào là hai biến cố xung khắc?
    \end{itemize}
\end{itemize}

\vspace{10pt}
\hrule
\vspace{10pt}

\begin{center}
    \textbf{\large TIẾT 2 (TIẾT 84 THEO PPCT | TUẦN 30)}\\[4pt]
    \textbf{\large BIẾN CỐ GIAO VÀ BIẾN CỐ XUNG KHẮC}
\end{center}

\subsection*{1. Hoạt động 1: Mở đầu (Khởi động - 6 phút)}

\begin{itemize}[leftmargin=1.2cm]
    \item[a)] \textbf{Mục tiêu:} Tạo tình huống có vấn đề để phân biệt sự khác nhau giữa từ "hoặc" (hợp) và từ "và" (giao).
    \item[b)] \textbf{Nội dung:}
    \begin{itemize}[leftmargin=0.6cm]
        \item Quay trở lại bài toán rút 1 lá bài tây ở Tiết 1:
        \begin{itemize}
            \item $A$: "Rút được lá Át".
            \item $B$: "Rút được lá chất Cơ".
        \end{itemize}
        \item Đặt câu hỏi: \textit{"Biến cố $D$: 'Lá bài rút ra VỪA LÀ LÁ ÁT VÀ VỪA MANG CHẤT CƠ' có bao nhiêu kết quả thuận lợi? Biến cố $D$ khác biến cố hợp $C$ ở Tiết 1 như thế nào?"}
    \end{itemize}

    \item[c)] \textbf{Sản phẩm:} Học sinh trả lời: Chỉ có duy nhất 1 lá bài thỏa mãn đồng thời cả hai điều kiện trên, đó chính là lá \textbf{"Át Cơ"}. Biến cố này đòi hỏi cả hai điều kiện cùng xảy ra đồng thời, hẹp hơn rất nhiều so với biến cố "hoặc" (có 16 lá).
    \item[d)] \textbf{Tổ chức thực hiện:} GV xác nhận câu trả lời rất chính xác và giới thiệu khái niệm Biến cố giao.
\end{itemize}

\subsection*{2. Hoạt động 2: Hình thành kiến thức mới (20 phút)}

\begin{itemize}[leftmargin=1.2cm]
    \item[a)] \textbf{Mục tiêu:}
    \begin{itemize}[leftmargin=0.6cm]
        \item Nắm vững định nghĩa biến cố giao của hai biến cố.
        \item Nắm vững khái niệm hai biến cố xung khắc (rời nhau, không thể cùng xảy ra).
        \item Thành thạo cách biểu diễn biến cố giao và xung khắc bằng sơ đồ Ven.
    \end{itemize}

    \item[b)] \textbf{Nội dung:}
    \begin{itemize}[leftmargin=0.6cm]
        \item \textbf{1. Định nghĩa biến cố giao:}
        \begin{itemize}
            \item Cho hai biến cố $A$ và $B$. Biến cố \textit{"Cả $A$ và $B$ cùng xảy ra"}, kí hiệu là $AB$ hoặc $A \cap B$, được gọi là \textbf{biến cố giao} của $A$ và $B$.
            \item Tập hợp các kết quả thuận lợi cho biến cố giao là giao của hai tập hợp $A$ và $B$:
            $$A \cap B = \{\omega \in \Omega \mid \omega \in A \text{ và } \omega \in B\}$$
            \item \textit{Sơ đồ Ven:} Miền biểu diễn $A \cap B$ là phần diện tích chung (giao nhau) của hai hình tròn $A$ và $B$.
        \end{itemize}

        \item \textbf{2. Định nghĩa biến cố xung khắc:}
        \begin{itemize}
            \item Hai biến cố $A$ và $B$ được gọi là \textbf{xung khắc} nếu chúng \textbf{không thể cùng xảy ra} trong một lần thực hiện phép thử.
            \item Nói cách khác, tập hợp các kết quả thuận lợi cho biến cố giao là tập hợp rỗng:
            $$A \cap B = \emptyset$$
            \item \textit{Sơ đồ Ven:} Hai hình tròn $A$ và $B$ hoàn toàn tách rời nhau, không có điểm chung nào.
        \end{itemize}
    \end{itemize}

    \item[c)] \textbf{Sản phẩm:} Học sinh ghi chép định nghĩa, vẽ sơ đồ Ven phân biệt $A \cap B$ và trường hợp xung khắc $A \cap B = \emptyset$.
    \item[d)] \textbf{Tổ chức thực hiện:}
\end{itemize}

\begin{pedabox}{Tiến trình tổ chức Hoạt động 2 (Kiến thức mới - Tiết 2)}
    \textbf{Bước 1: Chuyển giao nhiệm vụ}
    \begin{itemize}[leftmargin=0.6cm]
        \item GV vẽ bảng so sánh 2 cột: Biến cố hợp ($A \cup B$) và Biến cố giao ($A \cap B$).
        \item Yêu cầu học sinh thảo luận cặp đôi điền các từ khóa tương ứng: Phép toán tập hợp, Từ khóa ngôn ngữ, Ý nghĩa thực tế.
    \end{itemize}

    \textbf{Bước 2: Thực hiện nhiệm vụ có giàn giáo sư phạm}
    \begin{itemize}[leftmargin=0.6cm]
        \item HS trao đổi, quan sát hình ảnh sơ đồ Ven.
        \item GV hướng dẫn học sinh trung bình - yếu: Nhớ quy tắc "Hợp là gom hết - Giao là lấy chung".
    \end{itemize}

    \textbf{Bước 3: Báo cáo, thảo luận}
    \begin{itemize}[leftmargin=0.6cm]
        \item 1 cặp đôi báo cáo bảng so sánh; cả lớp thống nhất.
    \end{itemize}

    \textbf{Bước 4: Kết luận, chuẩn hóa}
    \begin{itemize}[leftmargin=0.6cm]
        \item GV chuẩn hóa bảng so sánh và định nghĩa hai biến cố xung khắc.
    \end{itemize}
\end{pedabox}

\begin{scaffoldbox}{Giàn giáo sư phạm cho học sinh trung bình - yếu (Scaffolding)}
    \textbf{Bảng phân biệt dứt khoát giữa Biến cố hợp và Biến cố giao:}
    \begin{center}
    \begin{tabular}{|l|c|c|}
    \hline
    \textbf{Tiêu chí} & \textbf{Biến cố hợp ($A \cup B$)} & \textbf{Biến cố giao ($A \cap B$)} \\ \hline
    \textbf{Từ khóa nhận diện} & \textbf{"HOẶC", "ít nhất một"} & \textbf{"VÀ", "đồng thời", "cả hai"} \\ \hline
    \textbf{Mô tả tập hợp} & $A \cup B$ (Gom tất cả lại) & $A \cap B$ (Chỉ lấy phần chung) \\ \hline
    \textbf{Trường hợp xung khắc} & $A \cap B = \emptyset$ (Hai biến cố không có phần chung) & Biến cố giao là biến cố không thể xảy ra! \\ \hline
    \end{tabular}
    \end{center}
\end{scaffoldbox}

\subsection*{3. Hoạt động 3: Luyện tập (14 phút)}

\begin{itemize}[leftmargin=1.2cm]
    \item[a)] \textbf{Mục tiêu:} Rèn luyện kỹ năng xác định biến cố giao và phân biệt hai biến cố có xung khắc hay không trong phép thử chọn học sinh.
    \item[b)] \textbf{Nội dung:} Thực hiện bài toán trên Phiếu học tập số 2:
    \begin{quote}
        \textbf{Ví dụ 2:} Lớp 11A trường THCS\&THPT Hồng Vân có 35 học sinh, trong đó có 20 bạn học giỏi môn Toán, 15 bạn học giỏi môn Vật lí và có 8 bạn học giỏi cả hai môn Toán và Vật lí. Chọn ngẫu nhiên một học sinh trong lớp. Gọi:
        \begin{itemize}
            \item $A$ là biến cố: "Học sinh được chọn học giỏi môn Toán".
            \item $B$ là biến cố: "Học sinh được chọn học giỏi môn Vật lí".
        \end{itemize}
        \begin{enumerate}
            \item Phát biểu bằng lời biến cố giao $AB$ và biến cố hợp $A \cup B$.
            \item Tính số học sinh thuận lợi cho biến cố $AB$ và biến cố $A \cup B$.
            \item Hai biến cố $A$ và $B$ có xung khắc với nhau không? Vì sao?
        \end{enumerate}
    \end{quote}

    \item[c)] \textbf{Sản phẩm (Lời giải chi tiết):}
    \begin{itemize}[leftmargin=0.6cm]
        \item \textit{Câu 1:}
        \begin{itemize}
            \item Biến cố giao $AB$: "Học sinh được chọn học giỏi cả hai môn Toán và Vật lí".
            \item Biến cố hợp $A \cup B$: "Học sinh được chọn học giỏi ít nhất một trong hai môn Toán hoặc Vật lí".
        \end{itemize}
        \item \textit{Câu 2:}
        \begin{itemize}
            \item Số kết quả thuận lợi cho $AB$ là số học sinh giỏi cả hai môn: $n(AB) = 8$ học sinh.
            \item Áp dụng công thức số phần tử của hợp hai tập hợp:
            $$n(A \cup B) = n(A) + n(B) - n(AB) = 20 + 15 - 8 = 27 \text{ học sinh}$$
        \end{itemize}
        \item \textit{Câu 3:}
        \begin{itemize}
            \item Hai biến cố $A$ và $B$ \textbf{không xung khắc} vì $n(AB) = 8 \ne 0$ (vẫn có 8 học sinh giỏi đồng thời cả hai môn, biến cố $AB$ hoàn toàn có thể xảy ra).
        \end{itemize}
    \end{itemize}

    \item[d)] \textbf{Tổ chức thực hiện:}
    \begin{itemize}[leftmargin=0.6cm]
        \item Học sinh làm việc cá nhân trong 7 phút. GV mời 1 HS lên bảng trình bày.
        \item Cả lớp đối chiếu kết quả; GV chốt công thức $n(A \cup B) = n(A) + n(B) - n(AB)$.
    \end{itemize}
\end{itemize}

\subsection*{4. Hoạt động 4: Vận dụng (5 phút)}

\begin{itemize}[leftmargin=1.2cm]
    \item[a)] \textbf{Mục tiêu:} Đặt tình huống dẫn dắt sang Tiết 3 về tính độc lập của các biến cố.
    \item[b)] \textbf{Nội dung & Câu hỏi suy ngẫm:}
    \begin{itemize}[leftmargin=0.6cm]
        \item Nếu hai bạn An và Bình ở hai phòng thi khác nhau cùng làm bài thi môn Toán. Biến cố $A$: "An đạt điểm 10", biến cố $B$: "Bình đạt điểm 10". Việc An đạt điểm 10 có làm ảnh hưởng gì đến khả năng đạt điểm 10 của Bình không?
        \item Chuẩn bị nội dung Tiết 3: Khái niệm biến cố độc lập.
    \end{itemize}
\end{itemize}

\vspace{10pt}
\hrule
\vspace{10pt}

\begin{center}
    \textbf{\large TIẾT 3 (TIẾT 85 THEO PPCT | TUẦN 30)}\\[4pt]
    \textbf{\large BIẾN CỐ ĐỘC LẬP VÀ ỨNG DỤNG TRONG MÔ HÌNH AI}
\end{center}

\subsection*{1. Hoạt động 1: Mở đầu (Khởi động - 6 phút)}

\begin{itemize}[leftmargin=1.2cm]
    \item[a)] \textbf{Mục tiêu:} Tiếp cận khái niệm biến cố độc lập qua tình huống gieo hai con xúc xắc riêng biệt.
    \item[b)] \textbf{Nội dung:}
    \begin{itemize}[leftmargin=0.6cm]
        \item Giáo viên phát cho mỗi nhóm bàn 2 con xúc xắc (1 màu xanh, 1 màu đỏ). Gieo đồng thời cả 2 con xúc xắc.
        \item Gọi $A$: "Con xúc xắc màu xanh xuất hiện mặt 6 chấm".
        \item Gọi $B$: "Con xúc xắc màu đỏ xuất hiện mặt 6 chấm".
        \item Đặt câu hỏi: \textit{"Nếu con xúc xắc màu xanh ra mặt 6 chấm thì có làm tăng hay giảm khả năng con xúc xắc màu đỏ ra mặt 6 chấm không?"}
    \end{itemize}

    \item[c)] \textbf{Sản phẩm:} Học sinh trả lời: Hoàn toàn không! Hai con xúc xắc hoàn toàn tách biệt nhau, con màu xanh ra mặt nào thì xác suất con màu đỏ ra mặt 6 chấm vẫn luôn luôn là $\dfrac{1}{6}$.
    \item[d)] \textbf{Tổ chức thực hiện:} GV dẫn dắt: Trong toán học, hai biến cố như vậy được gọi là "Hai biến cố độc lập".
\end{itemize}

\subsection*{2. Hoạt động 2: Hình thành kiến thức mới (20 phút)}

\begin{itemize}[leftmargin=1.2cm]
    \item[a)] \textbf{Mục tiêu:}
    \begin{itemize}[leftmargin=0.6cm]
        \item Nắm vững định nghĩa hai biến cố độc lập.
        \item Phân biệt rõ sự khác nhau giữa biến cố độc lập và biến cố xung khắc.
        \item Khám phá cách mô hình AI ứng dụng phân tích các sự kiện độc lập trong hệ thống gợi ý tìm kiếm (Mã 11.C5.2).
    \end{itemize}

    \item[b)] \textbf{Nội dung:}
    \begin{itemize}[leftmargin=0.6cm]
        \item \textbf{1. Định nghĩa hai biến cố độc lập:}
        \begin{itemize}
            \item Hai biến cố $A$ và $B$ được gọi là \textbf{độc lập} nếu sự xảy ra hay không xảy ra của biến cố này \textbf{không làm ảnh hưởng} đến xác suất xảy ra của biến cố kia.
            \item \textit{Ví dụ kinh điển về biến cố độc lập:}
            \begin{itemize}
                \item Gieo một đồng xu hai lần liên tiếp (kết quả lần 1 không ảnh hưởng đến lần 2).
                \item Hai xạ thủ cùng bắn độc lập vào một mục tiêu.
                \item Hai học sinh ở hai trường khác nhau làm bài thi tốt nghiệp.
            \end{itemize}
            \item \textit{Lưu ý quan trọng (Tránh sai lầm kinh điển):}
            \begin{itemize}
                \item \textbf{Biến cố xung khắc:} $A$ và $B$ không thể cùng xảy ra ($A \cap B = \emptyset$). Nếu $A$ xảy ra thì $B$ chắc chắn không xảy ra $\implies$ Chúng CÓ ảnh hưởng lẫn nhau, nên hai biến cố xung khắc (có xác suất dương) KHÔNG THỂ là hai biến cố độc lập!
                \item \textbf{Biến cố độc lập:} $A$ và $B$ VẪN CÓ THỂ cùng xảy ra đồng thời bình thường!
            \end{itemize}
        \end{itemize}
    \end{itemize}

    \item[c)] \textbf{Sản phẩm:} Học sinh ghi nhớ định nghĩa, hoàn thành bảng phân biệt giữa độc lập và xung khắc vào vở.
    \item[d)] \textbf{Tổ chức thực hiện:}
\end{itemize}

\begin{pedabox}{Tiến trình tổ chức Hoạt động 2 (Kiến thức mới - Tiết 3)}
    \textbf{Bước 1: Chuyển giao nhiệm vụ}
    \begin{itemize}[leftmargin=0.6cm]
        \item GV đưa ra câu hỏi trắc nghiệm thảo luận: \textit{"Nếu $A$ và $B$ là hai biến cố xung khắc thì chúng có độc lập với nhau không?"}
    \end{itemize}

    \textbf{Bước 2: Thực hiện nhiệm vụ có giàn giáo sư phạm}
    \begin{itemize}[leftmargin=0.6cm]
        \item Học sinh hay nhầm: Xung khắc là rời nhau nên tưởng là độc lập.
        \item GV phân tích: Giả sử $A$ là "Hôm nay trời mưa", $B$ là "Hôm nay trời không mưa". Hai biến cố này xung khắc. Nhưng nếu biết biến cố $A$ đã xảy ra (trời đang mưa), thì xác suất biến cố $B$ xảy ra ngay lập tức bằng $0$! Tức là sự xảy ra của $A$ đã làm thay đổi hoàn toàn xác suất của $B$. Do đó chúng phụ thuộc nhau, không hề độc lập!
    \end{itemize}

    \textbf{Bước 3: Báo cáo, thảo luận}
    \begin{itemize}[leftmargin=0.6cm]
        \item Cả lớp ồ lên thích thú khi nhận ra bản chất toán học sâu sắc này.
    \end{itemize}

    \textbf{Bước 4: Kết luận, chuẩn hóa}
    \begin{itemize}[leftmargin=0.6cm]
        \item GV đúc kết: Xung khắc là quan hệ loại trừ nhau; Độc lập là quan hệ không ảnh hưởng xác suất.
    \end{itemize}
\end{pedabox}

\begin{aibox}{Tích hợp Năng lực AI (Theo Quyết định 2422/QĐ-BGDĐT - Mã 11.C5.2)}
    \textbf{Ứng dụng giả định tính độc lập trong thuật toán gợi ý tìm kiếm AI:}
    \begin{itemize}[leftmargin=0.6cm]
        \item \textbf{Nguyên lý AI Recommender System (Bộ gợi ý YouTube, TikTok, Google):}
        \begin{itemize}
            \item Khi người dùng gõ từ khóa tìm kiếm (Ví dụ: "Toán 11" và "Học kỳ II"), mô hình AI xử lý ngôn ngữ tự nhiên (Naïve Bayes / Transformer) xem xét xác suất xuất hiện của các từ. Thuật toán phân loại Naïve Bayes thường đưa ra giả định đơn giản hóa rằng các từ xuất hiện độc lập với nhau trong văn bản để tính nhanh xác suất chủ đề của bài giảng.
            \item Trong hệ thống gợi ý xem video: Nếu biến cố $A$ ("Người dùng thích video Toán") và biến cố $B$ ("Người dùng thích video Bóng đá") được AI phân tích là độc lập, hệ thống sẽ đề xuất xen kẽ cả hai luồng video mà không để nội dung này chèn ép nội dung kia. Nếu phát hiện biến cố $A$ và $C$ ("Người dùng thích video Luyện thi ĐH") phụ thuộc chặt chẽ vào nhau, AI sẽ liên tục đẩy các video luyện thi vào trang chủ của người dùng.
        \end{itemize}
    \end{itemize}
\end{aibox}

\subsection*{3. Hoạt động 3: Luyện tập (14 phút)}

\begin{itemize}[leftmargin=1.2cm]
    \item[a)] \textbf{Mục tiêu:} Nhận diện chính xác các cặp biến cố độc lập và không độc lập trong các tình huống thực tiễn.
    \item[b)] \textbf{Nội dung:} Thực hiện bài toán trên Phiếu học tập số 3:
    \begin{quote}
        \textbf{Ví dụ 3:} Trong mỗi tình huống sau, hãy cho biết hai biến cố $A$ và $B$ có độc lập với nhau hay không? Giải thích vì sao.
        \begin{enumerate}[leftmargin=0.6cm]
            \item Gieo một đồng xu cân đối hai lần liên tiếp. Gọi $A$: "Lần thứ nhất xuất hiện mặt ngửa", $B$: "Lần thứ hai xuất hiện mặt ngửa".
            \item Một hộp có 3 viên bi xanh và 2 viên bi đỏ. Rút ngẫu nhiên lần thứ nhất 1 viên bi và \textbf{không trả lại vào hộp}, sau đó rút tiếp lần thứ hai 1 viên bi. Gọi $A$: "Lần thứ nhất rút được bi xanh", $B$: "Lần thứ hai rút được bi xanh".
            \item Hai bạn Nam và Việt mỗi người bắn một phát súng độc lập vào cùng một tấm bia. Gọi $A$: "Nam bắn trúng bia", $B$: "Việt bắn trúng bia".
        \end{enumerate}
    \end{quote}

    \item[c)] \textbf{Sản phẩm (Lời giải chi tiết):}
    \begin{itemize}[leftmargin=0.6cm]
        \item \textit{Trường hợp 1:} $A$ và $B$ là hai biến cố \textbf{độc lập}. Vì kết quả gieo của lần thứ nhất hoàn toàn không ảnh hưởng đến xác suất xuất hiện mặt ngửa của lần thứ hai (xác suất lần hai luôn là $\dfrac{1}{2}$).
        \item \textit{Trường hợp 2:} $A$ và $B$ \textbf{không độc lập} (phụ thuộc nhau). Vì sau lần thứ nhất rút bi mà \textit{không trả lại}:
        \begin{itemize}
            \item Nếu lần 1 rút được bi xanh (biến cố $A$ xảy ra), trong hộp còn lại 2 bi xanh và 2 bi đỏ $\implies$ Xác suất lần 2 rút được bi xanh là $\dfrac{2}{4} = \dfrac{1}{2}$.
            \item Nếu lần 1 rút được bi đỏ ($A$ không xảy ra), trong hộp còn lại 3 bi xanh và 1 bi đỏ $\implies$ Xác suất lần 2 rút được bi xanh là $\dfrac{3}{4}$.
            \item Như vậy việc biến cố $A$ xảy ra hay không đã làm thay đổi trực tiếp xác suất của biến cố $B$. Do đó $A$ và $B$ phụ thuộc nhau!
        \end{itemize}
        \item \textit{Trường hợp 3:} $A$ và $B$ là hai biến cố \textbf{độc lập}. Vì Nam và Việt bắn súng độc lập với nhau, kết quả bắn của Nam không làm ảnh hưởng đến khả năng bắn trúng bia của Việt.
    \end{itemize}

    \item[d)] \textbf{Tổ chức thực hiện:}
    \begin{itemize}[leftmargin=0.6cm]
        \item Cho học sinh thảo luận cặp đôi trong 6 phút. GV theo dõi, gợi ý học sinh phân tích số lượng bi còn lại trong hộp ở trường hợp 2.
        \item Mời đại diện 3 học sinh trả lời 3 trường hợp; cả lớp nhận xét, bổ sung.
    \end{itemize}
\end{itemize}

\subsection*{4. Hoạt động 4: Vận dụng (5 phút)}

\begin{itemize}[leftmargin=1.2cm]
    \item[a)] \textbf{Mục tiêu:} Khái quát hóa và dẫn dắt bài toán tính xác suất: Nếu $A$ và $B$ độc lập thì xác suất của biến cố giao $P(AB)$ được tính như thế nào?
    \item[b)] \textbf{Nội dung & Hướng dẫn về nhà:}
    \begin{itemize}[leftmargin=0.6cm]
        \item Ở trường hợp 3: Nếu xác suất Nam bắn trúng là $0{,}8$ và Việt bắn trúng là $0{,}7$. Dự đoán xem xác suất để cả hai cùng bắn trúng bia bằng bao nhiêu?
        \item Chuẩn bị cho bài học tiếp theo: Bài 29. Công thức cộng xác suất và Bài 30. Công thức nhân xác suất cho hai biến cố độc lập ($P(AB) = P(A) \cdot P(B)$).
    \end{itemize}
\end{itemize}

\section*{IV. PHỤ LỤC HỌC LIỆU BỔ TRỢ}

\subsection*{1. Phiếu học tập số 1 (Tiết 1): Biến cố hợp và Biểu diễn sơ đồ Ven}

\begin{pedabox}{PHIẾU HỌC TẬP SỐ 1: BÀI TẬP VỀ BIẾN CỐ HỢP}
    \textbf{Họ và tên học sinh:} \dotfill \textbf{Lớp:} 11\dots \textbf{Nhóm:} \dots
    
    \textbf{Nhiệm vụ 1: Hoàn thành định nghĩa}
    \begin{itemize}[leftmargin=0.6cm]
        \item Biến cố hợp của $A$ và $B$ là biến cố: \dotfill
        \item Kí hiệu là: \dotfill; Từ khóa ngôn ngữ là: \dotfill
    \end{itemize}

    \textbf{Nhiệm vụ 2: Bài tập rèn luyện}
    Gieo đồng thời hai con xúc xắc. Gọi $A$ là biến cố "Tổng số chấm xuất hiện bằng 7", $B$ là biến cố "Có ít nhất một con xúc xắc xuất hiện mặt 1 chấm".
    \begin{enumerate}
        \item Liệt kê các kết quả thuận lợi cho biến cố $A$: \dotfill
        \item Liệt kê các kết quả thuận lợi cho biến cố $B$: \dotfill
        \item Liệt kê các kết quả thuận lợi cho biến cố $A \cup B$: \dotfill
    \end{enumerate}
\end{pedabox}

\subsection*{2. Phiếu học tập số 2 (Tiết 2): Biến cố giao và Biến cố xung khắc}

\begin{pedabox}{PHIẾU HỌC TẬP SỐ 2: BIẾN CỐ GIAO VÀ BIẾN CỐ XUNG KHẮC}
    \textbf{Họ và tên học sinh:} \dotfill \textbf{Lớp:} 11\dots \textbf{Nhóm:} \dots

    \textbf{Nhiệm vụ 1: Điền đúng (Đ) hoặc Sai (S)}
    \begin{itemize}[leftmargin=0.6cm]
        \item Biến cố giao $AB$ xảy ra khi $A$ hoặc $B$ xảy ra. (\dots)
        \item Nếu hai biến cố $A$ và $B$ xung khắc thì chúng không có kết quả thuận lợi chung. (\dots)
        \item Nếu $A \cap B = \emptyset$ thì $A$ và $B$ là hai biến cố đối nhau. (\dots)
    \end{itemize}

    \textbf{Nhiệm vụ 2: Bài tập áp dụng}
    Rút ngẫu nhiên 1 thẻ từ hộp chứa 20 thẻ đánh số từ 1 đến 20. Gọi $A$: "Thẻ ghi số chia hết cho 3", $B$: "Thẻ ghi số chia hết cho 5".
    \begin{enumerate}
        \item Tìm các phần tử của biến cố giao $AB$: \dotfill
        \item Hai biến cố $A$ và $B$ có xung khắc không? Vì sao? \dotfill
    \end{enumerate}
\end{pedabox}

\subsection*{3. Phiếu học tập số 3 (Tiết 3): Biến cố độc lập}

\begin{pedabox}{PHIẾU HỌC TẬP SỐ 3: PHÂN BIỆT BIẾN CỐ ĐỘC LẬP}
    \textbf{Họ và tên học sinh:} \dotfill \textbf{Lớp:} 11\dots \textbf{Nhóm:} \dots

    \textbf{Xét tính độc lập của các cặp biến cố sau:}
    \begin{center}
    \begin{tabular}{|p{7.5cm}|c|p{5cm}|}
    \hline
    \textbf{Tình huống thực tế} & \textbf{Độc lập / Phụ thuộc} & \textbf{Giải thích ngắn gọn} \\ \hline
    1. Gieo 1 đồng xu và 1 con xúc xắc. $A$: "Đồng xu ngửa", $B$: "Xúc xắc 6 chấm". & \dots\dots\dots\dots & \dots\dots\dots\dots\dots\dots\dots\dots \\ \hline
    2. Rút 2 lá bài liên tiếp CÓ HOÀN LẠI. $A$: "Lần 1 rút lá K", $B$: "Lần 2 rút lá Q". & \dots\dots\dots\dots & \dots\dots\dots\dots\dots\dots\dots\dots \\ \hline
    3. Rút 2 lá bài liên tiếp KHÔNG HOÀN LẠI. $A$: "Lần 1 rút lá K", $B$: "Lần 2 rút lá Q". & \dots\dots\dots\dots & \dots\dots\dots\dots\dots\dots\dots\dots \\ \hline
    \end{tabular}
    \end{center}
\end{pedabox}

\subsection*{4. Rubric đánh giá năng lực học sinh trong bài học (4 mức độ)}

\begin{center}
\begin{tabularx}{\textwidth}{|>{\bfseries}p{3cm}|X|X|X|X|}
\hline
\textbf{Tiêu chí} & \textbf{Mức 1: Chưa đạt} & \textbf{Mức 2: Đạt} & \textbf{Mức 3: Khá} & \textbf{Mức 4: Tốt} \\ \hline
Nhận biết và biểu diễn biến cố hợp, giao & Nhầm lẫn giữa từ "hoặc" và "và", chưa vẽ được sơ đồ Ven. & Phân biệt được biến cố hợp và giao nhưng liệt kê phần tử còn trùng lặp. & Liệt kê chính xác kết quả thuận lợi, vẽ đúng sơ đồ Ven phân biệt hợp và giao. & Khái quát hóa xuất sắc mối liên hệ giữa lý thuyết tập hợp và xác suất, giải bài toán nhanh. \\ \hline
Phân biệt biến cố xung khắc và độc lập & Nhầm lẫn cho rằng biến cố xung khắc là độc lập. & Phân biệt được khái niệm nhưng giải thích lý do còn lúng túng. & Giải thích chuẩn xác tại sao hai biến cố xung khắc không thể là hai biến cố độc lập. & Phân tích sâu sắc bản chất xác suất có điều kiện, đưa ra các phản ví dụ thuyết phục. \\ \hline
Năng lực số (Mã 3.1NC1a) & Chưa vẽ được sơ đồ Ven trên máy tính hoặc điện thoại. & Vẽ được sơ đồ Ven đơn giản bằng công cụ Shapes trên Google Slides. & Thiết kế sơ đồ Ven số đẹp mắt, tô màu phân biệt các miền tập hợp rõ ràng. & Ứng dụng thành thạo phần mềm đồ họa, thiết kế hình trực quan sinh động chia sẻ cho lớp. \\ \hline
Nhận thức ứng dụng AI (Mã 11.C5.2) & Chưa hình dung được mối liên hệ giữa xác suất và AI. & Nêu được ví dụ AI gợi ý video dựa trên dữ liệu người dùng. & Hiểu được cách AI giả định các sự kiện độc lập trong thuật toán phân loại dữ liệu. & Phân tích sâu sắc cơ chế Recommender System, liên hệ với trải nghiệm số thực tế hàng ngày. \\ \hline
\end{tabularx}
\end{center}

\end{document}
